\documentclass[11pt]{article}
\usepackage[utf8]{inputenc}
\usepackage{amsmath,amssymb,amsthm,hyperref}
\newtheorem{theorem}{Theorem}
\newtheorem{lemma}{Lemma}
\newcommand{\R}{\mathbb R}
\newcommand{\E}{\mathbb E}
\newcommand{\Prb}{\mathbb P}
\newcommand{\A}{\mathcal A}
\title{Stationary parallel task allocation: a quota-bucket candidate}
\author{}
\date{20 October 2026}
\begin{document}
\maketitle
\noindent\textbf{Verification status.}
This is a mathematical derivation under review, not a completed Lean or Comparator
certificate. The setup, a causal measurable policy, its execution through the online
protocol, prefix statistics, uniform confidence, and confidence validity along the
actual trajectory are formalized separately. A proved reduction turns a
deterministic regret bound on the confidence event into an expected-regret bound
with an additive cost at most one. The schedule cost and the sum of phase costs
still require Lean proofs before the final theorem can be exported and audited
by Comparator. The Python implementation is a numerical research prototype.
Its tests do not certify the stochastic regret theorem below.

The Lean policy uses clipped intervals alone and updates them at phase boundaries,
after completing the planned quotas. The Python prototype additionally intersects
raw sub-Gaussian intervals and permits early phase completion. The argument below
applies to both variants; execution equivalence between the two implementations
has not been claimed.

\paragraph{Model and criterion.}
There are $n\ge1$ workers, $B\ge1$ tasks per round, and $K\ge1$ rounds.
For each worker, observations form an IID sequence; all sequences are mutually
independent. An observation $X_{i,j}$ is strictly positive almost surely, has
unknown mean $\mu_i\in(0,1]$, and satisfies
\[
 \E\exp(\lambda(X_{i,j}-\mu_i))\le\exp(\sigma^2\lambda^2/2),\qquad\lambda\in\R.
\]
The common upper bound $\sigma\ge0$ is known. Observations need not be at most one.
An action is $a\in\A=\{a\in\mathbb N^n:\sum_i a_i=B\}$; its feedback consists of
the next $a_i$ observations of worker $i$. There is no feedback before the action.
Write $\ell_v(a)=\max_i a_i v_i$,
$L(v)=\min_{a\in\A}\ell_v(a)$, $L^*=L(\mu)$, and
\[
 R_K=\sum_{k=1}^K(\ell_\mu(a_k)-L^*).
\]
The target is $\E R_K/K$, the expected average proxy regret, not realized makespan regret.
Every oracle below solves its finite integer minimization exactly. The heap
implementation selects the $B$ smallest slots $jv_i$, including a capacity
restriction when specified. Zero weights are permitted, and ties use worker index.

\paragraph{Explicit parameters.}
Put
\[
 d=B/n+n/B,\quad \bar\sigma=\max(1,\sigma),\quad
 t=\log\bigl(64nB^2K^3(1+\bar\sigma)\bigr),\quad
 H=1+\bar\sigma\sqrt{2t},\quad \chi=Ht,
\]
\[
 \beta=(1+\bar\sigma)e^{-t},\quad \varepsilon=\sqrt{d/K},\quad
 J_*=2+\left\lceil\log_2(162\chi B^2K)\right\rceil,\quad
 P_*=1+\left\lceil\log_2(1+B\sqrt K)\right\rceil.
\]
In particular, $t>1$, $\chi\ge1$, and $\beta\le\varepsilon/(16B)$.

\paragraph{Confidence intervals.}
Maintain nested intervals $[s_i,u_i]$, initially $[0,1]$. Store both the raw
observations and $Z_{i,j}=\min(X_{i,j},H)$. For $N_i>0$ and a previously valid
upper endpoint $v_i$, intersect the old interval and $[0,1]$ with the intervals
centered at the raw and clipped empirical means and having radii
\[
 r_i^{\rm raw}=\sigma\sqrt{2t/N_i},\qquad
 r_i^{\rm clip}=2\sqrt{v_i\chi/N_i}+2\chi/N_i+\beta.
\]
No update is made when $N_i=0$. Inside a phase, $v_i$ is frozen at that phase's
initial upper endpoint. Empty intersections trigger a feasible balanced fallback
for all remaining rounds. This also specifies the policy outside the confidence event.

\begin{lemma}[Uniform validity]
With probability at least $1-4nBK e^{-t}$, every interval update through round
$K$ contains its corresponding mean.
\end{lemma}
\begin{proof}
First consider a bounded random variable $0\le Z\le H$ with mean $\nu$.
For $|\lambda|H\le1$, the elementary inequality $e^x\le1+x+x^2$ for $|x|\le1$
and $Z^2\le HZ$ give
\[
 \E e^{\lambda(Z-\nu)}
 \le e^{-\lambda\nu}(1+\lambda\nu+\lambda^2H\nu)
 \le e^{\lambda^2H\nu}.
\]
Independence, exponential Markov, and the choice
$\lambda=\min\{\sqrt{t/(N\nu H)},1/H\}$ give, in both tails,
\[
 \Prb\left\{\left|N^{-1}\sum_{j=1}^N Z_j-\nu\right|
       >2\sqrt{\nu Ht/N}+2Ht/N\right\}\le2e^{-t}.
\]
For $\nu=0$, nonnegativity implies $Z=0$ almost surely and the assertion follows
directly. To check the optimization, if the square-root branch is selected,
the quadratic MGF term is $t$ and the square-root part of the threshold contributes
$2t$ to the negative exponent. On the other branch $N\nu/H\le t$, while the
linear threshold term contributes $2t$. Replacing $\lambda$ by $-\lambda$
proves the other tail.

For $Z=\min(X,H)$, let $\lambda=\sqrt{2t}/\bar\sigma$. Pointwise,
$(X-H)_+\le e^{\lambda(X-H)}/\lambda$. Since $\mu\le1$,
\[
 0\le\mu-\nu\le\lambda^{-1}
      e^{-\lambda(H-\mu)+\sigma^2\lambda^2/2}
 \le\frac{\bar\sigma}{\sqrt{2t}}e^{-t}\le\beta.
\]
Thus $\nu\le\mu\le v_i$ implies the clipped radius stated above. The raw
sub-Gaussian bound has failure probability $2e^{-t}$ for each prefix as well.
Take a union over workers, both interval constructions, and all prefix lengths
$1\le N\le BK$. On this single event, induction justifies every adaptive
choice of the previously valid $v_i$, without conditioning on adaptive counts.
\end{proof}

\paragraph{Policy.}
Set $\theta=B$. At a phase boundary compute the upper oracle value $V=L(u)$ on
the currently active set. Permanently remove each worker with $s_i>V$, using
strict inequality, and recompute the oracle values. Set
\[
 C_i=\begin{cases}\min\{B,\lfloor V/s_i\rfloor\},&s_i>0,\\B,&s_i=0,\end{cases}
 \qquad D_i=C_i(u_i-s_i)
\]
on the active set, and $C_i=D_i=0$ elsewhere. If every $D_i\le\theta/2$, halve
$\theta$ and repeat this check. Once $\theta\le\varepsilon$ or $V\le\varepsilon$,
use the upper oracle for all remaining rounds. Otherwise start a phase and freeze
$u,s,V,C,\theta$ for its allocation schedule. Give each active worker the quota
\[
 q_i=\left\lceil256\chi u_iC_i^2/\theta^2+32\chi C_i/\theta\right\rceil+1.
\]
Group workers by $2^j\le q_i<2^{j+1}$. In each group write $q=2^j$ and let $m$
be its size. Set $b=\min\{\min_i C_i,2q,\max(1,\lfloor B/m\rfloor)\}$.
Explore each worker for $2q$ new observations, in increments of at most $b$,
splitting the group into chunks of at most $\lfloor B/b\rfloor$ workers.
Process groups in increasing $j$, and chunks in worker order. This specifies
an exploration vector $e$ satisfying $e_i\le C_i$ and $\sum_i e_i\le B$ each round.

For every such vector, choose the upper-oracle residual allocation with capacities
$C_i-e_i$ and residual batch size $B-\sum_i e_i$. Play exploration plus residual.
Update intervals after receiving the entire batch, using the phase's frozen $u$
in its clipped radius. A phase may finish early when the updated geometry satisfies
$\max_i D_i\le\theta/2$. Only then halve $\theta$ and start the next phase.
At an exhausted horizon, record a pending phase as unfinished. Exhausting a quota
schedule without reaching its target triggers fallback, never a forced transition.
There is no positive-lower-bound warm-up in this variant.

The conservative capacity $\lfloor V/u_i\rfloor$ cannot replace $C_i$ in the
precision test. For $B=1$, $s=(0.2,0)$, $u=(0.2,1)$, and $\mu=(0.2,0.1)$,
these conservative capacities are $(1,0)$ and their corresponding uncertainties
are both zero. Nevertheless the upper oracle chooses worker one and has regret
$0.1$. Optimistic capacities correctly retain uncertainty on worker two.

\begin{lemma}[Oracle and interval geometry]
On the simultaneous confidence event, $L(s)\le L^*\le V$,
$0\le V-L(s)\le\max_i D_i$, and every true optimal allocation avoids permanently
eliminated workers. If $\max_i D_i\le\theta$, the upper oracle has regret at most
$\theta$, and every allocation satisfying $a_i\le C_i$ has regret at most $2\theta$.
For positive upper endpoints on $m$ active workers,
\[
 V\sum_{i\text{ active}}u_i^{-1}\le B+m-1\le B+n.
\]
Throughout execution, $V\le G:=\lceil B/n\rceil\le B/n+1$.
\end{lemma}
\begin{proof}
Monotonicity gives the sandwich. A lower-oracle minimizer $a^-$ satisfies
$a_i^-s_i\le L(s)\le V$, hence $a_i^-\le C_i$, including $s_i=0$. Therefore
$V\le\ell_u(a^-)\le L(s)+\max_iD_i$.
An eliminated worker has $\mu_i\ge s_i>V\ge L^*$ and cannot occur in any true
optimum. An upper-oracle minimizer itself never uses a worker with $s_i>V$,
so elimination does not increase $V$. Nested upper endpoints then make $V$
nonincreasing. Initially, balanced allocation at $u_i=1$ gives $V=G$.

For a capacity-feasible action, $a_is_i\le V$ and
$a_i(\mu_i-s_i)\le C_i(u_i-s_i)\le\theta$. Consequently its true loss is at most
$V+\theta$, and its regret at most $2\theta$. The upper oracle has loss at most $V$.

For the harmonic inequality let $c_i=\lceil V/u_i\rceil-1$. Every slot counted
in $c_i$ has cost strictly below $V$. If $\sum_i c_i\ge B$, trimming those counts
to total $B$ would give an allocation of loss strictly below the optimum $V$.
Thus $\sum_i c_i\le B-1$. Now sum $V/u_i\le c_i+1$.
The case $V=0$ is immediate; under the good event and $B\ge1$, $V>0$.
\end{proof}

\begin{lemma}[Feasible residual filling]
Let $b=\operatorname{RAS}(u,B)$, $b_i\le C_i$, and $e_i\le C_i$, $\sum_i e_i\le B$.
There is a residual $r$ with $r_i\le b_i$, $e_i+r_i\le C_i$, and
$\sum_i(e_i+r_i)=B$. In particular the capped residual oracle has loss at most $V$.
The resulting action has loss at most $V+\max_i e_i u_i$.
\end{lemma}
\begin{proof}
The residual capacities $\min\{b_i,C_i-e_i\}$ sum to at least
$\sum_i(b_i-e_i)=B-\sum_i e_i$. Trim them to the required residual batch.
This is a feasible residual of loss at most $\ell_u(b)=V$.
Finally $\max_i(e_i+r_i)\mu_i\le\max_i e_i u_i+\max_i r_i u_i$.
\end{proof}

\begin{lemma}[Quota-bucket scheduling]
For positive quotas and capacities let
$Q=\max_i q_i u_i$, $A=\max_i q_i/C_i$, and let $J$ be the number of nonempty
dyadic groups. The complete exploration schedule satisfies
\[
 N_{\rm round}\le4\sum_iq_i/B+2J(A+1),\qquad
 \sum_{\rm rounds}\max_i e_i u_i\le4Q(J+n/B).
\]
It collects between $q_i$ and $2q_i$ observations of each worker.
\end{lemma}
\begin{proof}
In a group $q\le q_i<2q$, hence $u_i\le Q/q$.
If $m\le B$, every worker fits in one chunk. There are $\lceil2q/b\rceil$ rounds,
at most $2q/\min_iC_i+2+2q/\lfloor B/m\rfloor
\le2A+2+4mq/B$. Total exploration load is at most
$(2q+b)Q/q\le4Q$.
If $m>B$, $b=1$ and the group needs $2q\lceil m/B\rceil\le4mq/B$ rounds, with
total exploration load at most $4Qm/B$. Sum over the groups and use
$mq\le\sum_{i\text{ in group}}q_i$.
\end{proof}

\begin{lemma}[Completion and phase costs]
Every phase starts with $\max_iD_i\le\theta$. If its complete quota schedule fits
within the horizon, the confidence event implies completion at its target.
With $G=\lceil B/n\rceil$, its regret, including any incomplete prefix, is at most
\[
 \begin{cases}
 20000\chi Jd,&\theta\ge V,\\
 21792\chi Jd/\theta+12J\theta,&\theta<V.
 \end{cases}
\]
\end{lemma}
\begin{proof}
The invariant holds initially since $C_i=B$, $u_i-s_i=1$ and $\theta=B$.
Nested intervals and decreasing $V$ make capacities and uncertainties nonincreasing.
After a completed phase it follows from the transition test.

After $q_i$ new samples, $N_i\ge q_i$ and the clipped interval's width is at most
\[
 4\sqrt{u_i\chi/N_i}+4\chi/N_i+2\beta
 \le\theta/(4C_i)+\theta/(8C_i)+\theta/(8C_i)=\theta/(2C_i).
\]
Updated capacities are at most the frozen $C_i$, proving completion.

Since $C_i u_i\le V+\theta$, the quotas imply
\[
 q_i u_i\le Q_0:=256\chi(V+\theta)^2/\theta^2+32\chi(V+\theta)/\theta+2,
\]
\[
 q_i/C_i\le A_0:=256\chi(V+\theta)/\theta^2+32\chi/\theta+2,
 \qquad \sum_iq_i\le Q_0(B+n)/V.
\]
The last inequality uses the harmonic oracle lemma. By the two regret bounds,
the phase costs at most both $2\theta N_{\rm round}$ and
$4Q_0(J+n/B)+(V-L^*)N_{\rm round}$, where $V-L^*\le\min(V,\theta)$.

For $\theta\ge V$, $Q_0\le1090\chi$ and $A_0\le544\chi/\theta+2$.
Substitution in the latter bound gives at most
$9808\chi J(1+n/B)+6JG\le20000\chi Jd$.
For $\theta<V$, $Q_0\le1090\chi V^2/\theta^2$ and
$A_0\le544\chi V/\theta^2+2$.
Substitution in the former bound gives at most
$10896\chi JV(1+n/B)/\theta+12J\theta$.
Use $V(1+n/B)\le(B/n+1)(1+n/B)=d+2\le2d$.

All bounds also hold for an unfinished phase: its frozen schedule extends
deterministically to a complete schedule, every summand of both bounding costs
is nonnegative, and no confidence assertion after the experimental horizon is used.
\end{proof}

\begin{theorem}[Candidate upper bound]
The policy described above satisfies
\[
 \frac{\E R_K}{K}\le
 \min\left\{B,\frac{100000\chi J_*(P_*+1)(d+\sqrt{dK})+1}{K}\right\}.
\]
\end{theorem}
\begin{proof}
Any started phase has $\theta>\varepsilon$. Its quotas are bounded by
$256\chi B^2/\varepsilon^2+32\chi B/\varepsilon+2
\le162\chi B^2K$, so $J\le J_*$. There are at most $P_*$ started phases.
The geometric precisions give $\sum_{\rm fine}1/\theta\le2/\varepsilon$.
Every fine precision is below $V\le G$, so $\sum_{\rm fine}\theta\le2G$.
Thus the exploration regret on the good event is at most
\[
 20000\chi J_*P_*d+43584\chi J_*d/\varepsilon+24J_*G.
\]
The final exploitation regret is at most $K\varepsilon$: stopping on precision
uses the interval lemma, and stopping on $V$ uses $\ell_\mu(\operatorname{RAS}(u))\le V$.
This also covers a run which starts directly in exploitation.
Use $G\le3d/2$, $d/\varepsilon=K\varepsilon=\sqrt{dK}$, and $\chi,J_*\ge1$.
On the complement event, every feasible action still has regret between zero
and $B$, so its contribution to expectation is at most
$BK\cdot4nBK e^{-t}\le1$. Integrability follows from $0\le R_K\le BK$.
\end{proof}

\paragraph{Lower-bound comparison and limits.}
For $n\ge2$ and $K\ge\max\{B/n,n/B\}$, the stationary specialization of
Theorem 13 of \href{https://arxiv.org/abs/2610.18265}{Begunov et al. (2026)}
gives $\Omega(\sqrt{nK/B})$ for $B<n$ and $\Omega(\sqrt{BK/n})$ for $B\ge n$.
Its positive observations are bounded by one, so they are a subclass of the
present model for a fixed suitable sub-Gaussian parameter. Since
$d$ is within a factor two of $\max\{B/n,n/B\}$, the candidate bound matches
these powers of $n,B,K$ up to logarithmic factors for fixed $\sigma>0$.
For smaller fixed positive $\sigma$, scale the bounded hard instances;
the constant in the lower bound changes. This comparison does not establish
optimal dependence on $\sigma$, a bound without logarithms, or a nonstationary result.
In the degenerate case $n=1$, regret is exactly zero.
The full stochastic theorem and its execution connection still require Lean
formalization and Comparator verification.
\end{document}
